% !TEX program = xelatex
% XeLaTeX, at least two passes. Preserve the supplied TUM Beamer template.
% PowerShell (from beamer/):
%   $env:TEXINPUTS='.;tum-templates-main/beamer;tum-templates-main/commons;tum-templates-main/commons/resources;'
%   xelatex -interaction=nonstopmode -halt-on-error tax_graph_pre.tex
% Source references are retained in Beamer notes; notes are not displayed.
\documentclass[english, sectionpage=false]{tumbeamer}
\usepackage{xeCJK}
\setCJKmainfont[BoldFont={Microsoft YaHei Bold}]{Microsoft YaHei}
\setCJKsansfont[BoldFont={Microsoft YaHei Bold}]{Microsoft YaHei}
\xeCJKsetup{CJKecglue={\hskip 0.15em plus 0.05em}}
\usepackage{booktabs}
\usepackage{tikz}
\usetikzlibrary{positioning, arrows.meta, calc}

\provideName{\theGroupName}{Self-defined Project Study}{Self-defined Project Study}
\provideName{\theSchoolName}{Cross-Border Tax Graph}{Cross-Border Tax Graph}
\TUMbeamersetup{
  title page = TUM tower,
  content page = TUM more space,
  headline = TUM threeliner,
  footline = TUM default,
  headline on = {every page},
  footline on = {content page},
}
\title{跨境税务决策支持}
\subtitle{项目价值、方法原理与下一阶段验证}
\author{Project Study}
\date{2026 年 10 月}
\footline{跨境税务决策支持~|~Project Study~|~2026}

\setbeamerfont{itemize/enumerate body}{size=\normalsize}
\setbeamerfont{itemize/enumerate subbody}{size=\small}
\setlength{\tabcolsep}{6pt}
\renewcommand{\arraystretch}{1.25}
\newcommand{\subnote}[1]{{\small\color{TUMGrayDark} #1}}
\newcommand{\sourcehint}[1]{{\scriptsize\color{TUMGrayDark} #1}}
\newcommand{\lead}[1]{{\large\bfseries\color{TUMBlue} #1}}
\newcommand{\stepnum}[1]{{\Large\bfseries\color{TUMBlue} #1}}
\setbeamertemplate{navigation symbols}{}

\newenvironment{specialframe}{\begingroup
  \setbeamertemplate{background canvas}[TUM flags]
  \setbeamertemplate{background}[TUM logo]
  \setbeamercolor{enumerate item}{fg=TUMWhite}
  \setbeamercolor{frametitle}{fg=TUMWhite}}{\endgroup}
\tikzset{
  box/.style={draw=TUMBlue, fill=TUMBlueLighter!30, rounded corners=2pt,
              align=center, font=\small, inner sep=5pt, minimum height=0.90cm},
  key/.style={box, fill=TUMBlue, text=white},
  res/.style={box, fill=TUMOrange!12, draw=TUMOrange},
  arr/.style={-{Latex[length=2mm]}, thick, TUMBlueDark},
  caption/.style={font=\small, align=center, text=TUMGrayDark},
}

\begin{document}
\maketitle
\note{[Sources]\par User-provided TUM Beamer class, tower illustration and logo assets in beamer/tum-templates-main. Original deck preserved in the working backup.\par[/Sources]}

\section{价值与交付}
\begin{frame}
  \frametitle{已跑通规则分析，下一步验证真实业务}
  \framesubtitle{当前阶段：可演示原型；范围：中国内地、香港、新加坡的公司}
  \vspace{0.2em}
  \lead{把集团税负、落地条件和原文依据放在一起比较}
  \vspace{0.8em}
  \begin{enumerate}
    \setlength{\itemsep}{0.65em}
    \item \textbf{已具备}\quad 结构化案例的方案比较、补问与已缴税评估。\\
      \subnote{输出税额区间、具体动作、办理期限和待确认条件。}
    \item \textbf{方法基础}\quad 固定规则计算，保留未知信息，并按保守税额选方案。\\
      \subnote{每项结果可追溯到规则、参数与官方条文。}
    \item \textbf{下一阶段}\quad 接通真实材料抽取，开展限定场景的专业复核。\\
      \subnote{先验证业务可用性，再决定范围扩展与模型优化。}
  \end{enumerate}
  \note{[Sources]\par tax_graph/run.py: main loads structured case JSON and calls planner; planner.py: plan, render; claims.py; llm/extract.py and check_extract.py. Current CLI does not call extraction.\par[/Sources]}
\end{frame}

\begin{frame}
  \frametitle{跨境方案需要同时回答三个问题}
  \framesubtitle{目标是支持可复核的决策，而不只是展示一个税率}
  \vspace{0.3em}
  \begin{tabular}{@{}p{0.32\textwidth}p{0.61\textwidth}@{}}
    \textcolor{TUMBlue}{\textbf{1}}\quad\textbf{整体成本} & 公司结构与交易安排怎样影响集团总税负？\\[0.9em]
    \textcolor{TUMBlue}{\textbf{2}}\quad\textbf{落地条件} & 哪些动作要先完成，哪些判断仍需取得认可？\\[0.9em]
    \textcolor{TUMBlue}{\textbf{3}}\quad\textbf{风险与期限} & 已缴税有哪些追回空间、补税风险和办理时限？
  \end{tabular}
  \par\vspace{1.0em}
  \lead{预期价值：统一比较口径，明确不确定性，减少遗漏}
  \par\vspace{0.35em}
  \subnote{节税收益与工作量改善仍需通过真实业务试点衡量。}
  \note{[Sources]\par Project objectives: v1_陆港新_企业.md. planner.py minimises group total under conservative completions. Actual ROI has not been measured; benefits are intended outcomes.\par[/Sources]}
\end{frame}

\begin{frame}
  \frametitle{信息充分给方案，信息不足给补问单}
  \framesubtitle{两种输出都说明依据和下一步动作}
  \begin{columns}[T]
    \begin{column}{0.47\textwidth}
      \lead{方案报告}\par\vspace{0.6em}
      \begin{itemize}
        \setlength{\itemsep}{0.55em}
        \item 建议结构与需要落实的动作
        \item 各法域及集团总税额区间
        \item 税项依据、期限与待认可条件
        \item 已缴税的追回空间与补税风险
      \end{itemize}
    \end{column}
    \begin{column}{0.47\textwidth}
      \lead{补问单}\par\vspace{0.6em}
      \begin{itemize}
        \setlength{\itemsep}{0.55em}
        \item 缺少哪项数据或确认
        \item 通常从哪些材料取得
        \item 答复为什么会影响方案
        \item 哪些缺口必须优先补齐
      \end{itemize}
    \end{column}
  \end{columns}
  \vspace{0.7em}
  \subnote{答复进入案例后重新计算；无法界定税额的候选，先补齐必要信息。}
  \note{[Sources]\par planner.py: render; run.py writes reports to out/. Answers currently supplied in a separate answers JSON. Paid-tax analysis: claims.py.\par[/Sources]}
\end{frame}

\begin{frame}
  \frametitle{同一案例贯通方案比较与已缴税评估}
  \framesubtitle{当前入口是结构化案例，计算始终使用已编码的规则}
  \centering\vspace{0.8em}
  {\small\textcolor{TUMBlue}{\textbf{共同依据：官方法条、参数、规则与生效日期}}}
  \par\vspace{0.7em}
  \begin{tikzpicture}[x=1cm,y=1cm]
    \draw[arr] (0.96,0)--(1.69,0);
    \draw[arr] (3.61,0)--(4.34,0);
    \draw[arr] (6.26,0)--(6.99,0);
    \draw[arr] (8.91,0)--(9.64,0);
    \draw[arr,TUMOrange] (7.95,-0.55)--(7.95,-1.4)--(0,-1.4)--(0,-0.55);
    \node[box,text width=1.60cm,minimum height=1.1cm] at (0,0) {集团事实\\\scriptsize 公司、交易};
    \node[box,text width=1.60cm,minimum height=1.1cm] at (2.65,0) {候选结构\\\scriptsize 模板与安排};
    \node[key,text width=1.60cm,minimum height=1.1cm] at (5.30,0) {税额区间\\\scriptsize 规则与条件};
    \node[key,text width=1.60cm,minimum height=1.1cm] at (7.95,0) {比较补问\\\scriptsize 保守端优先};
    \node[res,text width=1.60cm,minimum height=1.1cm] at (10.60,0) {方案报告\\\scriptsize 动作与期限};
    \node[caption,text=TUMOrange] at (3.95,-1.65) {补齐事实或确认后，重新计算};
  \end{tikzpicture}
  \par\vspace{0.9em}
  \subnote{已缴税按同一组计算区间评估：可追回金额、补税风险与处理路线。}
  \note{[Sources]\par run.py; scenarios.py: candidates and graph; engine.py: bounds; planner.py: selection and questions; claims.py. Diagram simplifies existing workflow. Retrieval and extraction are not live mandatory calculator steps.\par[/Sources]}
\end{frame}

\section{理论与方法}
\begin{frame}
  \frametitle{图模型把持股关系与交易分开表示}
  \framesubtitle{同一结构既有控制关系，也有不同方向的资金流}
  \centering
  \begin{tikzpicture}[x=1cm,y=1cm]
    \draw[arr] (1.3,0.30)--(2.7,0.30);
    \draw[arr] (5.3,0.30)--(6.7,0.30);
    \draw[arr,TUMOrange] (2.7,-0.30)--(1.3,-0.30);
    \draw[arr,TUMOrange] (6.7,-0.30)--(5.3,-0.30);
    \node[box,text width=2.2cm,minimum height=1.3cm] at (0,0) {母公司\\\textbf{新加坡}};
    \node[box,text width=2.2cm,minimum height=1.3cm] at (4,0) {控股公司\\\textbf{香港}};
    \node[key,text width=2.2cm,minimum height=1.3cm] at (8,0) {运营公司\\\textbf{内地}};
    \node[caption,text=TUMBlue] at (2,0.95) {持股 100\%};
    \node[caption,text=TUMBlue] at (6,0.95) {持股 100\%};
    \node[caption,text=TUMOrange] at (2,-0.95) {股息};
    \node[caption,text=TUMOrange] at (6,-0.95) {股息};
  \end{tikzpicture}
  \par\vspace{0.7em}
  \begin{itemize}
    \setlength{\itemsep}{0.3em}
    \item \textbf{持股边}\quad 记录比例、取得日期与成本，派生比例和持有期。
    \item \textbf{交易流}\quad 明确收入类型、收付双方、金额与发生日期。
    \item 各法域分别计算交易税项，再合并集团承担的总税额。
  \end{itemize}
  \sourcehint{三层控股结构为教学示意，不是客户方案；候选搜索限于当前模板与范围。}
  \note{[Sources]\par model.py: Entity, HoldEdge, FlowEdge and Graph; scenarios.py: Template and Case; facts.py: facts_for and holding derivations; engine.py. Teaching structure follows synthetic CASE_A in cases.py: SG parent, HK holder, CN operating company, each wholly owned; dividend flows in the opposite direction. No client facts or claimed business results.\par[/Sources]}
\end{frame}

\begin{frame}
  \frametitle{规则把法律条件转成可复算的判断}
  \framesubtitle{先编码适用范围、条件与效果，再对案例字段执行比较}
  \centering
  \begin{tikzpicture}[x=1cm,y=1cm]
    \draw[arr] (2.05,0.75)--(3.0,0.75)--(3.0,0)--(3.8,0);
    \draw[arr] (2.05,-0.75)--(3.0,-0.75)--(3.0,0);
    \draw[arr] (5.6,0)--(6.5,0);
    \node[box,text width=2.65cm] at (0.55,0.75) {条件 A\\\textbf{居民证明已取得}};
    \node[box,text width=2.65cm] at (0.55,-0.75) {条件 B\\\textbf{持股门槛已满足}};
    \node[key,text width=1.35cm] at (4.7,0) {A 且 B\\\scriptsize AND};
    \node[res,text width=2.95cm] at (8.15,0) {两项均为真\\\textbf{该组条件满足}};
  \end{tikzpicture}
  \par\vspace{0.6em}
  \begin{itemize}
    \setlength{\itemsep}{0.35em}
    \item 数字门槛从参数表读取，并核对原文引用与生效日期。
    \item 例外、反滥用与抵免等规则按预定关系继续处理。
  \end{itemize}
  \sourcehint{条件组合示意，省略其他法定条件；不代表已满足完整优惠资格。}
  \note{[Sources]\par rules/model.py: Leaf, And, Or, Not, RuleObj; rules/eval.py: cond and stages; rules/compile.py. 规则层.md sections 1, 3-5. Conditions form a teaching example, not a complete treaty rule.\par[/Sources]}
\end{frame}

\begin{frame}
  \frametitle{未知信息保留下来，并进入税额区间}
  \framesubtitle{三值逻辑：条件可以是真、假或未知}
  \begin{columns}[T]
    \begin{column}{0.47\textwidth}
      \lead{“且”的判断示例}\par\vspace{0.6em}
      \begin{tabular}{@{}lll@{}}
        \toprule 条件 A & 条件 B & 结果\\
        \midrule
        真 & 真 & \textcolor{TUMBlue}{\textbf{真}}\\[0.25em]
        假 & 未知 & \textbf{假}\\[0.25em]
        真 & 未知 & \textcolor{TUMOrange}{\textbf{未知}}\\
        \bottomrule
      \end{tabular}
      \par\vspace{0.6em}
      \subnote{一项已确定不满足时，不必为该条件组继续补齐另一项。}
    \end{column}
    \begin{column}{0.47\textwidth}
      \lead{税额保留两端}\par\vspace{0.8em}
      \begin{tikzpicture}[x=1cm,y=1cm]
        \draw[very thick,TUMBlueLight] (0,0)--(5.0,0);
        \draw[thick,TUMBlue] (0,-0.12)--(0,0.12);
        \draw[thick,TUMOrange] (5,-0.12)--(5,0.12);
        \node[caption,anchor=north,text=TUMBlue] at (0,-0.22) {有利端\\\scriptsize 最好情形};
        \node[caption,anchor=north,text=TUMOrange] at (5,-0.22) {不利端\\\scriptsize 保守情形};
        \node[caption] at (2.5,0.42) {待定条件影响区间宽度};
      \end{tikzpicture}
      \par\vspace{1.1em}
      \subnote{缺数据、未确认的动作与机关裁量分开记录；补齐后再收窄区间。}
    \end{column}
  \end{columns}
  \note{[Sources]\par rules/tri.py: AND, OR, NOT, Tri, Need; engine.py: bounds and non-monotone groups; planner.py. 伪代码.md sections 2-5. Original presentation identifies strong three-valued logic with Libkin (2016). Interval is conceptual, not a confidence interval or forecast.\par[/Sources]}
\end{frame}

\begin{frame}
  \frametitle{按保守税额选方案，减少乐观依赖}
  \framesubtitle{稳健决策：比较各方案在不利情形下的集团税负}
  \centering
  \begin{tikzpicture}[x=0.16cm,y=0.90cm]
    \draw[arr,TUMGray] (0,-1.2)--(69,-1.2);
    \foreach \x in {0,10,20,30,40,50,60} {
      \draw[TUMGray] (\x,-1.25)--(\x,-1.15);
      \node[font=\scriptsize,text=TUMGrayDark,anchor=north] at (\x,-1.35) {\x};
    }
    \draw[line width=5pt,TUMOrange!35] (10,0.85)--(60,0.85);
    \draw[line width=5pt,TUMBlueLighter] (25,-0.35)--(40,-0.35);
    \fill[TUMOrange] (10,0.85) circle (2.5pt);
    \fill[TUMOrange] (60,0.85) circle (2.5pt);
    \fill[TUMBlue] (25,-0.35) circle (2.5pt);
    \fill[TUMBlue] (40,-0.35) circle (2.5pt);
    \node[font=\small\bfseries,anchor=east] at (-2,0.85) {方案 A};
    \node[font=\small\bfseries,anchor=east] at (-2,-0.35) {方案 B};
    \node[caption,anchor=south] at (10,1.03) {最好 10};
    \node[caption,anchor=south,text=TUMOrange] at (60,1.03) {保守 60};
    \node[caption,anchor=south] at (25,-0.15) {最好 25};
    \node[caption,anchor=south,text=TUMBlue] at (40,-0.15) {保守 40};
    \node[font=\scriptsize,text=TUMGrayDark,anchor=north west] at (65,-1.35) {税负};
  \end{tikzpicture}
  \par\vspace{0.6em}
  \lead{本例选择 B：保守税额 40，低于 A 的 60}
  \par\vspace{0.35em}
  \subnote{即使 A 的最好情形更低，也不把尚未确认的有利结果作为承诺。}
  \par\vspace{0.3em}
  \sourcehint{教学示例：数字与单位为示意，不是客户测算。方法：$\Gamma$-maximin。}
  \note{[Sources]\par planner.py: _key orders ceiling, floor, pending discretion and entity count. 优化方案.md and original deck cite Troffaes (2007) for Gamma-maximin. A=[10,60], B=[25,40] are hypothetical teaching intervals; min(60,40)=40.\par[/Sources]}
\end{frame}

\begin{frame}
  \frametitle{优先补齐会改变决策的信息}
  \framesubtitle{信息价值：看答复是否可能改变方案或降低保守税额}
  \centering
  \begin{tikzpicture}[x=1cm,y=1cm]
    \draw[arr] (2.35,0)--(3.10,0);
    \draw[arr] (6.1,0)--(6.85,0);
    \node[box,text width=2.15cm] at (1.1,0) {当前选择 B\\\textbf{保守税额 40}};
    \node[res,text width=2.70cm] at (4.6,0) {关键问题 Q\\\scriptsize 有利答复获确认时\\A 的保守税额降至 30};
    \node[key,text width=2.50cm] at (8.25,0) {重新选择 A\\\textbf{保守税额 30}};
  \end{tikzpicture}
  \par\vspace{0.8em}
  \begin{itemize}
    \setlength{\itemsep}{0.35em}
    \item \textbf{必须问}\quad 缺失使税额无法界定的数据。
    \item \textbf{优先问}\quad 能影响方案、动作可行性或关键裁定的信息。
    \item 多项条件共同影响结果时，列出需要一起确认的事项。
  \end{itemize}
  \sourcehint{延续上一页教学示例；潜在影响为 40 − 30 = 10，不是按成功概率计算的预期收益。}
  \note{[Sources]\par planner.py: plan computes forced favourable ceilings; _joint_pairs checks jointly valuable questions. Unknown answers remain unknown; rulings must be confirmed. Diagram is hypothetical, not an actual result.\par[/Sources]}
\end{frame}

\begin{frame}
  \frametitle{AI 辅助找依据，材料抽取仍待接通}
  \framesubtitle{混合检索综合词面、含义与条款关系}
  \centering
  \begin{tikzpicture}[x=1cm,y=1cm]
    \draw[arr] (1.0,0)--(1.65,0);
    \draw[arr] (1.65,0)--(1.65,1.2)--(2.2,1.2);
    \draw[arr] (1.65,0)--(2.2,0);
    \draw[arr] (1.65,0)--(1.65,-1.2)--(2.2,-1.2);
    \draw[arr] (5.7,1.2)--(6.2,1.2)--(6.2,0)--(6.65,0);
    \draw[arr] (5.7,0)--(6.65,0);
    \draw[arr] (5.7,-1.2)--(6.2,-1.2)--(6.2,0);
    \node[box,text width=1.0cm] at (0.35,0) {问题};
    \node[box,text width=3.15cm] at (3.95,1.2) {编号与关键词\\\scriptsize 精确匹配 + BM25};
    \node[key,text width=3.15cm] at (3.95,0) {含义相近的文本\\\scriptsize BGE-M3 向量模型};
    \node[box,text width=3.15cm] at (3.95,-1.2) {显式条款关系\\\scriptsize 引用与图邻居};
    \node[res,text width=2.50cm] at (8.10,0) {融合排名\\\textbf{相关条文与原文}};
  \end{tikzpicture}
  \par\vspace{0.2em}
  \begin{itemize}
    \setlength{\itemsep}{0.25em}
    \item \textbf{已配置}\quad BGE-M3 检索；税额与方案选择由固定规则执行。
    \item \textbf{待接通}\quad 真实 LLM 从合同、财报中抽取事实与证据。
  \end{itemize}
  \sourcehint{AI 可在开发阶段辅助编码规则；规则仍需原文核对和专业复核。}
  \note{[Sources]\par official/README.md sections 7, 9, 10; official/tools/search.py; official/_index/dense_model.txt: BAAI/bge-m3. RRF combines exact, BM25, dense and graph channels. https://huggingface.co/BAAI/bge-m3; Cormack et al. (2009), Reciprocal Rank Fusion. llm/extract.py; check_extract.py uses scripted responder. run.py does not call extraction or retrieval as its calculator.\par[/Sources]}
\end{frame}

\begin{frame}
  \frametitle{用独立答案与可推导关系核验结果}
  \framesubtitle{验证计算正确性，也检查边界和依据是否完整}
  \centering
  \begin{tikzpicture}[x=1cm,y=1cm]
    \draw[arr] (2.55,0.65)--(3.25,0.65)--(3.25,0)--(3.9,0);
    \draw[arr] (2.55,-0.65)--(3.25,-0.65)--(3.25,0);
    \draw[arr] (5.7,0)--(6.5,0);
    \node[box,text width=2.50cm] at (1.15,0.65) {官方算例\\\textbf{独立期望结果}};
    \node[box,text width=2.50cm] at (1.15,-0.65) {同一输入\\\textbf{系统计算结果}};
    \node[key,text width=1.30cm] at (4.8,0) {比较};
    \node[res,text width=2.75cm] at (8.0,0) {一致：保留证据\\差异：定位并修正};
  \end{tikzpicture}
  \par\vspace{0.55em}
  \begin{itemize}
    \setlength{\itemsep}{0.30em}
    \item \textbf{关系与边界}\quad 金额缩放、门槛两侧、条件反转及抵免上限。
    \item \textbf{依据核对}\quad 引语、条款、时点，以及事实是否附来源。
    \item \textbf{业务确认}\quad 真实材料与专业复核，验证方案是否适用、可落实。
  \end{itemize}
  \note{[Sources]\par 测试系统.md sections 0, 2, 3, 8; check_examples.py, check_sg_s45.py, battery.py, mutation.py, verify/grounding.py. Some checks have text-derived expectations; diagram shows official-example subset. Saved snapshots detect change, not independent correctness.\par[/Sources]}
\end{frame}

\section{进展与下一步}
\begin{frame}
  \frametitle{已具备演示基础，业务证据仍需补齐}
  \framesubtitle{资料与规则快照：2026 年 10 月 14 日}
  \begin{columns}[T]
    \begin{column}{0.30\textwidth}
      \centering{\fontsize{26}{29}\selectfont\bfseries\color{TUMBlue}3}\par
      \small 当前法域\\\subnote{内地、香港、新加坡}
    \end{column}
    \begin{column}{0.30\textwidth}
      \centering{\fontsize{26}{29}\selectfont\bfseries\color{TUMBlue}472}\par
      \small 入库官方文件\\\subnote{含原件分片}
    \end{column}
    \begin{column}{0.30\textwidth}
      \centering{\fontsize{26}{29}\selectfont\bfseries\color{TUMBlue}250}\par
      \small 当前编译规则\\\subnote{有引用与生效区间}
    \end{column}
  \end{columns}
  \vspace{0.8em}
  \begin{tabular}{@{}p{0.14\textwidth}p{0.79\textwidth}@{}}
    \textcolor{TUMBlue}{\textbf{已有}} & 结构化分析、补问与方案报告、自动核验、部分公开案例比对。\\[0.5em]
    \textcolor{TUMOrange}{\textbf{待验证}} & 抽取准确性、方案可落实性、机关认可条件，以及新增规则覆盖。
  \end{tabular}
  \par\vspace{0.6em}
  \subnote{既有测试支持原型验证；实际业务收益与效率改善将在试点中衡量。}
  \note{[Sources]\par official/STATUS.md: 472 documents including 70 fragments, 551 parameters. Import of engine.RULES on 2026-10-14: 250 compiled rules. model.py: CN, HK, SG. 案例采集.md, 测试系统.md, out/*.grounding.md: disclosed/semisynthetic cases. No deployment readiness or customer ROI claimed. Captured snapshot.\par[/Sources]}
\end{frame}

\begin{frame}
  \frametitle{下一阶段先验证一个限定业务场景}
  \framesubtitle{建议路径：真实案例 → 可靠事实 → 可复核方案}
  \begin{enumerate}
    \setlength{\itemsep}{0.8em}
    \item \textbf{选定试点与专业基准}\\
      \subnote{优先从现有股息或服务费场景选择，取得完整材料与专业判断。}
    \item \textbf{接通真实材料抽取}\\
      \subnote{验证字段准确性、原文证据和缺失识别；再送入现有分析链路。}
    \item \textbf{复核方案、解释差异}\\
      \subnote{检查税额、待定条件、办理路径与可落实性，并补齐回归测试。}
  \end{enumerate}
  \vspace{0.6em}
  \lead{用试点证据决定范围扩展与模型优化的优先级}
  \note{[Sources]\par Proposed next-stage sequence, not an approved schedule or resource commitment. Existing templates: scenarios.py. Extraction: llm/extract.py. Expert review: v1_陆港新_企业.md. Real/private cases: 案例采集.md. No budget, durations, named owners or acceptance thresholds invented.\par[/Sources]}
\end{frame}

\begin{specialframe}
\begin{frame}
  \vspace*{0.10\textheight}
  {\color{TUMWhite}\LARGE\bfseries 明确试点范围与专业支持}\par\vspace{0.8em}
  \begin{minipage}{0.62\textwidth}
    \color{TUMWhite}
    \textbf{建议本轮确认}\par\vspace{0.5em}
    \begin{enumerate}
      \setlength{\itemsep}{0.4em}
      \item 优先验证哪一类跨境交易
      \item 案例材料与三地复核的提供方
      \item 试点通过所需的验收证据
    \end{enumerate}
    \vspace{0.7em}
    \textbf{试点交付}\par\vspace{0.2em}
    {\small 可追溯方案、未决条件清单、验证结果}
  \end{minipage}
  \note{[Sources]\par Proposed discussion items and deliverables, not approved decisions or commitments. Supplied TUM flag background retained.\par[/Sources]}
\end{frame}
\end{specialframe}

\appendix
\section{补充材料}
\begin{frame}
  \frametitle{补充：核心组件与证据链}
  \framesubtitle{将资料、规则、事实、计算与核验分开维护}
  \small
  \renewcommand{\arraystretch}{1.08}
  \begin{tabular}{@{}p{0.16\textwidth}p{0.46\textwidth}p{0.30\textwidth}@{}}
    \toprule
    \textbf{组件} & \textbf{职责} & \textbf{主要位置}\\
    \midrule
    知识库 & 条文、参数、检索、引文及时效 & \texttt{official/}\\
    规则层 & 范围、条件、效果与编译校验 & \texttt{rules/}\\
    图与事实 & 实体、持股、交易与事实派生 & \texttt{model / facts}\\
    计算引擎 & 税项、区间、跨流交互与未决项 & \texttt{engine.py}\\
    规划器 & 候选比较、信息价值与补问 & \texttt{planner.py}\\
    已缴税 & 追回、补税风险、时限与路线 & \texttt{claims / ex\_post}\\
    核验与输出 & 算例、溯源、先例因素与报告 & \texttt{verify / run}\\
    抽取接口 & 字段与原文；真实模型待接通 & \texttt{llm/extract.py}\\
    \bottomrule
  \end{tabular}
  \par\vspace{0.5em}
  \subnote{先例分析用于辅助依据核验；当前默认规划入口不自动据此确认裁量。}
  \note{[Sources]\par Project modules inspected. precedent.py is used in verify/grounding.py, not planner.py or default CLI. Compiler checks provenance and references, not substantive correctness of all legal encoding.\par[/Sources]}
\end{frame}

\begin{frame}
  \frametitle{补充：验证记录与适用边界}
  \framesubtitle{现有报告按保存批次解释，不等同于当前代码全量复跑}
  \small
  \begin{tabular}{@{}p{0.18\textwidth}p{0.29\textwidth}p{0.44\textwidth}@{}}
    \toprule
    \textbf{证据} & \textbf{已保存记录} & \textbf{如何解读}\\
    \midrule
    案例电池 & 1,058 例，0 缺陷 & 合成测试；该报告评估 237 条规则。\\[0.15em]
    触发覆盖 & 3 条规则未触发 & 新增规则及未触发规则仍需补测。\\[0.15em]
    规划器对照 & 78 对补全后方案一致 & 历史批次，验证补问后的结果一致性。\\[0.15em]
    原文溯源 & 条款、引语与时点核对 & 还需检查案例事实是否附可靠来源。\\[0.15em]
    专业复核 & 待开展限定场景试点 & 计算通过不等于业务适用或机关认可。\\
    \bottomrule
  \end{tabular}
  \par\vspace{0.6em}
  \sourcehint{资料：out/battery\_report.md、out/check\_planner.txt、out/*.grounding.md。}
  \par\vspace{0.4em}
  \subnote{方法依据：三值逻辑、稳健决策与信息价值、混合检索、独立核验及关系测试。}
  \note{[Sources]\par battery_report.md read on 2026-10-14: 1058 clean cases; rules never fired 3 of 237: hk.property.trader, hk.property.trader.none, hk.stamp.property.seller. check_planner.txt dated 2026-10-08: 78 same, 0 different, 0 defects. Current rule import: 250. Separate snapshots, not fresh full validation. gold_A.grounding.md shows missing factual provenance despite no failed legal-grounding checks. 测试系统.md and 案例采集.md.\par[/Sources]}
\end{frame}
\end{document}
